% =====================================================================
%  Chapter 07 — Layer 4
% =====================================================================
\chapter{Layer 4: Diagnostics, and What They Could and Could Not Decide}
\label{ch:layer4}

Layer~4 is the largest layer by module count and the one with the most
complicated relationship to the mathematics. Each module was built to measure
a specific quantity appearing in a specific labelled statement of the proof
document. Some of those measurements were decisive about the mechanism they
probed. One large family of them was, as the programme eventually proved, incapable
of deciding anything at all --- and that finding, Wall~6, is as much a result
of this chapter as any number in it.

\section{The diagnostic inventory}
\label{sec:layer4-inventory}

\begin{table}[ht]
\centering
\footnotesize
\begin{tabular}{@{}l l p{0.42\textwidth}@{}}
\toprule
Module & Session & What it measures, and for which statement \\
\midrule
\texttt{delta\_extractor.py} & \Sn{10}
  & \(\delta(t)\) and Sobolev norms of \(\theta\); the Claim-A exponent. \\
\texttt{cz\_integrability.py} & \Sn{10}
  & \(\norm{S_\theta}_{L^p}\) for \(p\in\{1,1.1,1.25,1.5,2\}\); how far the
    Calder\'on--Zygmund gain reaches. \\
\texttt{lps\_monitor\_3d.py} & \Sn{10}
  & \(\norm{u}_{L^p}\), enstrophy, palinstrophy, Prodi--Serrin margin. \\
\texttt{mean\_morrey\_diagnostic.py} & \Sn{23}, \Sn{25}
  & \(r^{-1+\alpha}\abs{\langle u\rangle_{B_r}}^3\); the mean-Morrey gap of
    \S14 of the proof document. \\
\texttt{gamma\_adversarial\_search.py} & \Sn{27}
  & \(\Gamma(r)\) under adversarial data; the \S15 gradient-concentration
    condition. \\
\texttt{geometric\_disorder.py} & \Sn{30}
  & \(M\), \(\rstar\), \(A_{\mathrm{loc}}\), \(\sigma=A_{\mathrm{loc}}/M^{3/2}\),
    \(R_{\log}\). \\
\texttt{alignment\_bridge.py} & \Sn{32}, \Sn{33}
  & The bridge ratio \(\norm{\nabla\ehat}_{L^\infty}\rstar/\sqrt{\sstar}\);
    \texttt{conj:bridge} (line~5466). \\
\texttt{k\_correlation.py} & \Sn{34}
  & \(c_K\), the decorrelation constant of \texttt{conj:K-refined}
    (line~5917). \\
\texttt{kcorr\_reynolds\_sweep.py} & \Sn{35}--\Sn{37}
  & \(c_K\) against viscosity and resolution. \\
\texttt{depletion\_diagnostic.py} & \Sn{36}
  & \(\alpha(x^*)/M\); the depletion hypothesis H2. \\
\texttt{enstrophy\_exponent.py} & \Sn{37}
  & \(p\) in \(\Omega\sim M^p\); the Regime~A/B dichotomy. \\
\texttt{aposteriori\_verification.py} & \Sn{39}
  & Residual-based Gr\"onwall margin (a diagnostic, never a proof). \\
\texttt{band\_concentration.py} & \Sn{42}, \Sn{43}
  & Band/bulk concentration ratios; \texttt{target:band-absorption}
    (line~8435). \\
\texttt{scale\_diagnostics.py} & \Sn{45}
  & \(\abs{E}/\abs{Q(\rstar)}\), \(c_K^{\mathrm{band}}\), and the
    \(\Lfac\)-dependence demanded by \texttt{prop:s44-linfty-route}
    (line~9209). \\
\texttt{continuation\_criterion.py} & \Sn{16}
  & \(A_{\min}\), \(\mu_{\min}\) against the Phase-0 criterion. \\
\bottomrule
\end{tabular}
\caption[The Layer-4 diagnostic inventory]{The Layer-4 modules, in the order
they were built, with the statement each targets. The chronology matters: the
modules from \Sn{32} onward all serve the alignment programme of Part~IV, and
\S\ref{sec:wall6} explains why that whole family could not deliver what it was
asked for.}
\label{tab:layer4-inventory}
\end{table}

\section{The early diagnostics: Claim~A made visible}
\label{sec:early-diagnostics}

The first three modules exist to make the Claim-A exponent measurable. Their
verdict, from \Sn{10}:

\begin{center}
\small
\begin{tabular}{@{}l l r l@{}}
\toprule
\texttt{exp\_id} & IC & key metric & verdict \\
\midrule
\texttt{EXP-L4-CZ-TG-032} & Taylor--Green & \(\varepsilon_{\mathrm{CZ}}=1.00\); Claim~A on \(100\%\) of steps & \textsf{PASS} \\
\texttt{EXP-L4-CZ-SHEAR-032} & shear & \(\varepsilon_{\mathrm{CZ}}=1.00\); Claim~A on \(100\%\) of steps & \textsf{PASS} \\
\texttt{EXP-L4-LPS-TG-032} & Taylor--Green & \(Z_{\max}=4.14\!\times\!10^{-1}\), \(\norm{u}_{L^6}\le0.639\) & \textsf{PASS} \\
\texttt{EXP-L4-LPS-SHEAR-032} & shear, \(\varepsilon=0\) & \(Z_{\max}=1.703\), \(\norm{u}_{L^6}\le0.998\) & \textsf{PASS} \\
\bottomrule
\end{tabular}
\end{center}

\medskip
\(\varepsilon_{\mathrm{CZ}}=1.00\) means the dissipation \(S_\theta\) was
controlled all the way up to \(L^2\), i.e.\ one full unit of integrability
beyond the \(L^1\) the energy inequality gives. The five \(\delta\)-extractor
rows \texttt{EXP-L4-DELTA-EPS-TG-032-eps\{1.0,0.1,0.01,0.001,0.0\}} all return
\(\delta_{\min}=\delta_{\max}=\delta_{\mathrm{mean}}=1.50\), positive from
\(t=0.05\) onwards, at every \(\varepsilon\) including zero. This is the same
\(\varepsilon\)-independence seen in Chapters~\ref{ch:routes} and
\ref{ch:layer3}, measured a third way.

\section{A lesson about sampling}
\label{sec:sampling-lesson}

The mean-Morrey diagnostic produced the programme's cleanest methodological
warning, and it is worth a section because the failure mode is not specific to
this programme.

The quantity is a supremum: \(\sup_r r^{-1+\alpha}\abs{\langle
u\rangle_{B_r}}^3\), a worst-case bound over balls. The first implementation
sampled eight ball centres on a fixed grid. For Taylor--Green and shear-layer
fields --- which are built from sines and cosines --- those centres land on
symmetry nodes where the mean vanishes identically by construction. The
measured maximum came out at \(7.70\times10^{-47}\).

\begin{center}
\small
\begin{tabular}{@{}l l l r r@{}}
\toprule
\texttt{exp\_id} & IC & sampling & \(M_\alpha^{\max}\) & \(\gamma\) \\
\midrule
\texttt{EXP-L4-MM-TG-001} & TG & deterministic, 8 fixed & \(7.70\!\times\!10^{-47}\) & \(+0.090\) \\
\texttt{EXP-L4-MM-SH-001} & shear & deterministic, 8 fixed & \(2.72\!\times\!10^{-47}\) & \(+0.055\) \\
\texttt{EXP-L4-MM-TG-002} & TG & random, 64, seed 42 & \(\mathbf{11.03}\) & \(+0.139\) \\
\texttt{EXP-L4-MM-SH-002} & shear & random, 64, seed 42 & \(\mathbf{14.82}\) & \(+0.384\) \\
\bottomrule
\end{tabular}
\end{center}

\medskip
Forty-seven orders of magnitude of pure sampling bias. The scaling exponent
\(\gamma\) --- the quantity the diagnostic actually exists to estimate --- is
positive in all four cases and so the qualitative conclusion survives, but the
magnitudes in the first two rows are meaningless. Both sets of rows are
retained in the database, which is why the comparison can be made at all.

\begin{record}
\textbf{The rule this produced.} When the quantity is a supremum over a
family, sample the family randomly with a logged seed, not on a lattice that
shares the field's symmetries. The programme's later diagnostics
(\texttt{scale\_diagnostics.py}, \texttt{band\_concentration.py}) sample at the
vorticity maximum rather than on a grid, for the same reason.
\end{record}

\section{A negative result, correctly read}
\label{sec:gamma-negative}

\texttt{EXP-L4-GAM-ADV-001} is the programme's most-quoted single negative
row, and its \texttt{verdict} field is not \textsf{PASS} or \textsf{FAIL} but
a sentence:

\begin{record}
\texttt{NOT SUPPRESSED: NS evolution does not increase Gamma exponent
(0.000 -> -3.263). Possible gap in Claim A conjecture for adversarial ICs.}
\end{record}

The run constructs a divergence-free bump concentrated at scale \(r\),
evolves it, and measures whether \(\Gamma(r)\) is suppressed. It is not: the
fitted exponent goes to \(-3.263\).

The \Sn{27} log's reading is the correct one and is worth reproducing,
because it is a model of how to handle a result that appears to contradict
one's own conjecture. The evolution in that run is \emph{free viscous decay},
\(\partial_tu=\nu\Delta u\), with no nonlinear term. The mechanism the
conjecture appeals to --- the scale-recursive redistribution of energy ---
lives entirely in the nonlinearity that was omitted. So the run does not
falsify the conjecture; it establishes that linear diffusion alone does not
suppress gradient concentration, which nobody had claimed. The stated next
step was to rerun the adversarial data through the full spectral solver.

That rerun was never performed under this identifier. The row stands in the
database, with its sentence, unresolved.

\section{The alignment diagnostics}
\label{sec:alignment-diagnostics}

From \Sn{32} the diagnostics serve the alignment programme, whose central
open estimate is the Bridge Lemma (\texttt{conj:bridge}, line~5466): does
\emph{averaged} coherence of the vorticity direction field upgrade to
\emph{pointwise} coherence? The measurable proxy is the bridge ratio
\(R=\norm{\nabla\ehat}_{L^\infty,\mathrm{high}}\,\rstar/\sqrt{\sstar}\), which
the lemma predicts should stay \(O(1)\).

\begin{table}[ht]
\centering
\small
\begin{tabular}{@{}l l r r r r@{}}
\toprule
\texttt{exp\_id} & IC & \(N\) & \(R_{\max}\) & \(R_{\mathrm{med}}\) &
\(\sstar\) range \\
\midrule
\texttt{EXP-L4-BRIDGE-TG-001} & TG & \(32\) & \(0.982\) & \(0.682\) & \([0.62,\,58.5]\) \\
\texttt{EXP-L4-BRIDGE-SH-001} & shear & \(32\) & \(0.719\) & \(0.639\) & \([12.6,\,35.8]\) \\
\texttt{EXP-L4-BRIDGE-TG-002} & TG & \(64\) & \(1.264\) & \(0.764\) & \([0.47,\,177.9]\) \\
\texttt{EXP-L4-BRIDGE-SH-002} & shear & \(64\) & \(0.763\) & \(0.587\) & \([4.01,\,128.1]\) \\
\texttt{EXP-L4-BRIDGE-ADV-001} & anti-par.\ tubes & \(64\) & \(1.364\) & \(1.059\) & \([3.18,\,9.38]\) \\
\texttt{EXP-L4-BRIDGE-ADV-002} & tubes, \(\nu=5\!\times\!10^{-4}\) & \(64\) & \(1.521\) & \(0.995\) & \([1.28,\,2.77]\) \\
\bottomrule
\end{tabular}
\caption[The bridge-ratio measurements]{All six bridge-ratio runs. The lemma's
own ceiling in the conjectured form is a constant \(C_B\); the measured ratio
never exceeded \(1.53\) in any run, including a \(4\times\) refinement and an
escalated-amplitude, lower-viscosity adversarial run.}
\label{tab:bridge}
\end{table}

The scientifically interesting row is \texttt{EXP-L4-BRIDGE-ADV-001}. The
anti-parallel tube configuration --- a direction-reversal interface, precisely
where a Bridge Lemma counterexample would live --- runs at \(R\) in a tight
band \([0.92,1.36]\) for the whole simulation, \emph{nearly saturating} the
inequality without ever violating it. Pointwise and averaged coherence move in
near-exact lockstep. Meanwhile Taylor--Green stresses the other axis: \(M\)
grows by \(10.5\times\) and \(R\) still stays below \(1.27\).

The caveat is the one that governs this entire family of measurements and is
stated in \S\ref{sec:wall6}: these are globally regular flows at moderate
Reynolds number, and the lemma is a statement about the near-singular regime.

\section{The decorrelation constant, and four sessions of it}
\label{sec:ck}

\texttt{conj:K-refined} (line~5917) asks for a decorrelation constant
\(c_K\lesssim C/\Lfac\) near a putative Type-I singularity. The programme
measured \(c_K\) in \Sn{34}, \Sn{35}, \Sn{42}--\Sn{43} and \Sn{45}.

\begin{center}
\small
\begin{tabular}{@{}l l r r r l@{}}
\toprule
\texttt{exp\_id} & IC & \(\nu\) & \(N\) & \(c_K\) median & note \\
\midrule
\texttt{EXP-L4-KCORR-TG-001} & TG & \(10^{-3}\) & \(64\) & \(1.063\) & \\
\texttt{EXP-L4-KCORR-SH-001} & shear & \(10^{-3}\) & \(64\) & \(1.935\) & worst split residual \\
\texttt{EXP-L4-KCORR-ADV-001} & tubes & \(10^{-3}\) & \(64\) & \(1.156\) & well resolved \\
\texttt{EXP-L4-KCORR-NU-TG-002} & TG & \(5\!\times\!10^{-4}\) & \(64\) & \(1.021\) & tail frac.\ \(0.33\) \\
\texttt{EXP-L4-KCORR-NU-TG-003} & TG & \(2.5\!\times\!10^{-4}\) & \(64\) & \(1.020\) & tail frac.\ \(0.59\) \\
\texttt{EXP-L4-KCORR-NU-ADV-003} & tubes & \(2.5\!\times\!10^{-4}\) & \(64\) & \(1.113\) & tail frac.\ \(6\!\times\!10^{-5}\) \\
\texttt{EXP-L4-KCORR-NU-TG-101} & TG & \(10^{-3}\) & \(\mathbf{128}\) & \(\mathbf{2.235}\) & resolved \\
\texttt{EXP-L4-KCORR-NU-TG-102} & TG & \(5\!\times\!10^{-4}\) & \(\mathbf{128}\) & \(\mathbf{2.238}\) & resolved \\
\texttt{EXP-L4-KCORR-NU-TG-103} & TG & \(2.5\!\times\!10^{-4}\) & \(\mathbf{128}\) & \(\mathbf{2.239}\) & resolved \\
\bottomrule
\end{tabular}
\end{center}

\medskip
Two findings, one of them a correction of the other.

The \Sn{35} headline was that \(c_K\) does not grow with Reynolds number: over
a fourfold decrease in viscosity it sat at \(1.02\)--\(1.16\) with regression
slopes of \(-0.031\) against \(\log(1/\nu)\) for both Taylor--Green and the
adversarial tubes. That removed the one numerical failure mode --- strong
positive correlation, \(c_K\gg1\) --- that could have defeated the absorption
argument outright.

The \Sn{36}--\Sn{37} correction is that the \(N=64\) Taylor--Green numbers were
\emph{resolution-suppressed by roughly a factor of two}. The tail fraction of
the spectrum at \(N=64\) ranged from \(0.07\) to \(0.59\); at \(N=128\) it fell
to \(1.3\times10^{-5}\)--\(2.6\times10^{-5}\), and the measured \(c_K\) rose
from \(\approx1.05\) to \(\approx2.24\). That is systematic spectral pile-up at
the grid cutoff, not noise. Taylor--Green sits noticeably \emph{above} the
adversarial series, not beside it as the under-resolved comparison had
suggested. A further probe, \texttt{EXP-L4-KCORR-NU-TG-104}, showed the tail
fraction climbing back to \(0.0082\) late in the Taylor--Green cascade even at
\(N=128\): resolution adequacy is time-window-dependent, not a property of
\((N,\nu)\) alone. That caveat was flagged and never resolved.

\section{The measurement the closure route actually needs}
\label{sec:l-dependence}

\Sn{44} identified two closure routes for the remaining core term, and its
\texttt{prop:s44-linfty-route} (line~9209) requires not boundedness of the
correlation constant but \emph{decay}: \(c_K\le C_1/\Lfac\). \Sn{45} built
\texttt{scale\_diagnostics.py} to measure exactly that, computing every
quantity inside the ball \(B_{\rstar}(x^*)\) rather than over the whole box,
and regressing \(\log c_K\) on \(\log\Lfac\).

\begin{figure}[ht]
\centering
\begin{tikzpicture}
\begin{axis}[
  width=0.82\textwidth, height=6.6cm,
  xlabel={\(\Lfac\) (median over samples)},
  ylabel={\(c_K\) (median over samples)},
  xmin=1.2, xmax=3.7, ymin=0, ymax=2.0,
  grid=both, grid style={gray!18},
  legend pos=north east, legend cell align=left,
  legend style={font=\scriptsize, draw=gray!40},
  tick label style={font=\scriptsize}, label style={font=\small}
]
\addplot[only marks, mark=*, color=openblue] coordinates {
  (3.3189,1.1013) (3.1710,1.0000) (1.7460,1.1812)
  (3.3805,1.1003) (3.4809,1.0000) (2.0984,1.2048)
};
\addlegendentry{TG / shear / tubes, two viscosities}
\addplot[only marks, mark=square*, color=provedgreen] coordinates {
  (1.4111,1.0749) (1.3767,1.0911) (1.3886,1.0716)
  (1.4371,1.0257) (1.6064,0.9819) (2.0603,0.9178)
};
\addlegendentry{jet strain sweep}
\addplot[dashed, thick, impossiblered, domain=1.2:3.7, samples=40]
  {1.5/x};
\addlegendentry{what closure needs: \(c_K\propto1/\Lfac\)}
\addplot[thick, gray, domain=1.2:3.7, samples=2] {1.08};
\addlegendentry{what was measured: flat}
\end{axis}
\end{tikzpicture}
\caption[The \(\Lfac\)-decay test]{The measurement the closure route requires,
and what came back. A clean \(C/\Lfac\) law (dashed) would fall by a factor of
about \(2.5\) across the measured range. The twelve runs instead sit pinned
between \(0.92\) and \(1.21\). The fitted slope of \(\log c_K\) on
\(\log\Lfac\) is \(-0.196\) (\(r^2=0.05\)) on the first pool and \(+0.088\)
(\(r^2=0.02\)) across all twelve, against the \(-1\) required. The low
\(r^2\) is not a failure of the fit --- it is the finding.}
\label{fig:l-decay}
\end{figure}

The same session made the first measurement of the volume hypotheses at scale
\(\rstar\). The bulk hypothesis (NB-E) is comfortably satisfied:
\(\abs{E}/\abs{Q(\rstar)}\) ranges over \(0.45\)--\(0.99\) with
\(\Lfac\approx1.7\)--\(3.5\), giving a product of \(1.4\)--\(2.6\). The band
hypothesis (NB) is not: \(\abs{\Band}/\abs{Q(\rstar)}\) ranges over
\(0.00\)--\(0.18\), and in one run --- shear at \(\nu=10^{-3}\) --- it is
\emph{exactly zero}: the transition band \(\{M/8\le\abs{\omega}\le M/4\}\) was
empty inside \(B_{\rstar}\), with all the vorticity there lying above
\(M/4\). Where that happens the band question is locally vacuous and the
volume route is unavailable precisely there.

So both of the two routes \Sn{44} laid out came back unsupported: the first
needs a decay that cannot be detected, and the second needs a hypothesis that
measurably fails.

\section{The jet sweep: the aligned corner, probed on purpose}
\label{sec:jet-sweep}

The one place the alignment programme predicts something must happen is the
\(\sstar\to0\) corner: a singularity, if there is one, must approach it.
\Sn{45} built a bipolar swirling jet and swept the strain-to-swirl ratio at
fixed energy, driving the flow deliberately toward alignment.

\begin{center}
\small
\begin{tabular}{@{}l r r r r r@{}}
\toprule
\texttt{exp\_id} & strain ratio & \(\abs{E}/\abs{Q}\) &
\(\abs{\Band}/\abs{Q}\) & \(c_K\) & \(\sstar\) \\
\midrule
\texttt{EXP-L4-JET000-064} & \(0.00\) & \(0.632\) & \(0.141\) & \(1.075\) & \(11.72\) \\
\texttt{EXP-L4-JET025-064} & \(0.25\) & \(0.609\) & \(0.159\) & \(1.091\) & \(14.91\) \\
\texttt{EXP-L4-JET050-064} & \(0.50\) & \(0.553\) & \(0.170\) & \(1.072\) & \(8.05\) \\
\texttt{EXP-L4-JET100-064} & \(1.00\) & \(0.459\) & \(0.159\) & \(1.026\) & \(6.71\) \\
\texttt{EXP-L4-JET200-064} & \(2.00\) & \(0.300\) & \(0.125\) & \(0.982\) & \(\mathbf{4.31}\) \\
\texttt{EXP-L4-JET400-064} & \(4.00\) & \(0.405\) & \(0.136\) & \(0.918\) & \(21.54\) \\
\bottomrule
\end{tabular}
\end{center}

\medskip
\(\sstar\) falls from \(11.72\) at pure swirl to a minimum of \(4.31\) at
strain ratio \(2\), then jumps to \(21.54\) --- consistent with the jet
destabilising once strain dominates swirl. The descent confirms the physical
picture that motivated the experiment. And \(c_K\) stays between \(0.92\) and
\(1.09\) across the entire sweep, \emph{including at the \(\sstar\)
minimum}. The aligned corner was probed on purpose, and the correlation
constants did not degrade there: mild positive evidence for their robustness,
and a null result for the hope that this geometry would expose adversarial
structure.

A resolution caveat is flagged in the source and repeated here: the ball
\(B_{\rstar}\) at \(N=64\) contains only \(81\) grid cells for Taylor--Green.
Shear has \(147\)--\(179\); the adversarial runs about \(1650\)--\(1750\); the
jet runs \(1300\)--\(4000\). A Taylor-Green-only reading of these numbers would
not be trustworthy; the better-resolved runs tell the same story, which is why
the conclusions are stated at all.

\section{The band conjecture, and a correction to a correction}
\label{sec:band}

\texttt{target:band-absorption} (line~8435) asks that the coupling density on
the transition band not be concentrated there relative to the bulk. \Sn{42}
built the diagnostic and ran it at \(N=64\); \Sn{43} repeated at \(N=128\), at
a cost of about \(2.7\) hours per run.

\begin{center}
\small
\begin{tabular}{@{}l l r r r r r@{}}
\toprule
\texttt{exp\_id} & IC & \(N\) & \(\mathrm{conc}_Y\) med & \(\mathrm{conc}_Y\) max &
\(\mathrm{conc}_D\) med & \(\mathrm{conc}_D\) max \\
\midrule
\texttt{EXP-L4-BAND-TG-001} & TG & \(64\) & \(0.460\) & \(26.57\) & \(1.308\) & \(63.15\) \\
\texttt{EXP-L4-BAND-SH-001} & shear & \(64\) & \(0.061\) & \(2.301\) & \(6.505\) & \(192.77\) \\
\texttt{EXP-L4-BAND-ADV-001} & tubes & \(64\) & \(1.526\) & \(9.810\) & \(2.890\) & \(5.184\) \\
\texttt{EXP-L4-BAND-TG-101} & TG & \(128\) & \(0.405\) & \(135.95\) & \(1.299\) & \(214.44\) \\
\texttt{EXP-L4-BAND-SH-101} & shear & \(128\) & \(0.378\) & \(0.950\) & \(3.916\) & \(75.01\) \\
\bottomrule
\end{tabular}
\end{center}

\medskip
The medians are resolution-robust: Taylor--Green's \(\mathrm{conc}_D\) median
moves by \(0.7\%\) across a \(2\times\) refinement. The maxima are not, and
they move in \emph{opposite directions} --- Taylor--Green's grew fivefold,
shear's shrank. A statistic that moves in opposite directions under refinement
on two initial conditions is not converged, and nothing can rest on it.

\Sn{42} had read the spikes as ``likely resolution artefacts''.
\Sn{43} records that this was \emph{half wrong} and replaces it: the spikes
grew under refinement, so they are not resolution artefacts in the sense
supposed. Splitting each run's timeseries at \(t=3\) settles what they are ---
in every run, at both resolutions, on all three IC families, \emph{every
maximum occurs at \(t<3\)}, in the early transient, while the high-vorticity
set has not yet been developed by nonlinear mixing and band statistics over
few unrepresentative points are volatile. Taylor--Green at \(N=128\), the run
with the largest spike in the study at \(135.95\), never exceeds \(0.63\)
anywhere in the bulk; its bulk maxima agree with the \(N=64\) run to within
\(7\)--\(9\%\).

The bulk regime is the one the conjecture concerns. There the density is tame,
stable across initial conditions and across refinement. The conjecture is
numerically supported --- to exactly the standard the programme applied to
\(c_K\), which is to say ``supported, not proved'', and which
\S\ref{sec:wall6} is about to qualify much more sharply.

\section{The a posteriori diagnostic, and a specification error}
\label{sec:aposteriori}

\texttt{aposteriori\_verification.py} implements a verification-\emph{style}
diagnostic in the shape of the Chernyshenko--Constantin--Robinson--Titi
technique. Its own documentation is emphatic that it does not produce a proof:
a rigorous computer-assisted argument requires validated interval arithmetic on
the residual and on every norm entering the Gr\"onwall estimate, and everything
here is ordinary floating point with an underived constant. The module never
returns a certificate, only a margin.

The first two runs returned margins of \(5.51\times10^{19}\) and
\(2.02\times10^{14}\), reported \textsf{PARTIAL}. The addendum to the session
log states what happened: the runs were \emph{vacuous}, and the fault was in
the specification. Accumulating \(\exp\int_0^TA\,ds\) over a twenty-time-unit
window with \(A\sim O(1\text{--}3.5)\) gives \(e^{70}\approx10^{30}\), so the
margin could not have been met by any residual whatsoever. The underlying
trajectories were perfectly well behaved throughout: the \(H^{-1}\) residual
peaked at \(4.0\times10^{-3}\) for Taylor--Green and stayed below \(0.074\)
for shear.

The windowed rerun fixed it, resetting the error at each window start:

\begin{center}
\small
\begin{tabular}{@{}l r r r r@{}}
\toprule
IC & window \(W\) & windows & fraction satisfied & max Gr\"onwall factor \\
\midrule
TG & \(0.5\) & \(26\) & \(0.423\) & \(5.11\) \\
TG & \(1.0\) & \(13\) & \(0.154\) & \(28.6\) \\
TG & \(2.0\) & \(7\)  & \(0.000\) & \(920\) \\
shear & \(0.5\) & \(23\) & \(\mathbf{0.957}\) & \(3.37\) \\
shear & \(1.0\) & \(12\) & \(0.917\) & \(9.81\) \\
shear & \(2.0\) & \(6\)  & \(0.833\) & \(79.2\) \\
\bottomrule
\end{tabular}
\end{center}

\medskip
The Gr\"onwall factor fell from \(\sim10^{12}\) to \(O(3\text{--}5)\), so the
comparisons became informative --- but neither trajectory reaches a clean pass
at any window length, and both experiments remain \textsf{PARTIAL}. The
shear-versus-Taylor--Green difference is real but reflects the size of the
numerical residual against an arbitrary threshold, not either flow's proximity
to a singularity. One suspected artefact was flagged and not fixed: shear's
very first window fails at all three lengths, most likely because a one-sided
time difference at step zero gives a poor estimate of \(\partial_tu\).

\section{Wall 6: why most of this chapter could not decide anything}
\label{sec:wall6}

The programme measured \(c_K\) in four separate sessions. In \Sn{47} it proved
that none of those measurements could ever have been evidence, and the proof
does not depend on any number any run produced.

\begin{record}
\textbf{Wall 6} (\texttt{subsec:wall6-numerics}). \texttt{conj:K-refined} and
\texttt{conj:lc-s44} are \emph{conditional} statements: each asserts a bound
\emph{near a putative Type-I singularity}, in the limit \(\Lfac\to\infty\). A
solver that remains globally regular is never in that regime. So neither the
conjecture nor its negation predicts anything about the numbers such a solver
produces --- both predict the same thing, namely nothing --- and a measurement
that both sides of a disjunction predict cannot separate them.

\smallskip
\texttt{prop:regime-A-impossible} (line~6514) removes the one reading under
which a regular run could still have served as a proxy. It proves that at a
genuine blow-up time \(\limsup_{t\to T^-}\Omega(t)=\infty\); a resolved
simulation is a bounded-enstrophy object throughout, so the regime the
conjectures speak about is not merely one the runs did not visit --- it is not
of the type a resolved run represents.

\smallskip
Classification: \IMPOSSIBLE{} as a test of the conjectures.
\end{record}

The consequence is a prohibition, and it is stated in the document: measuring
\(c_K\), \(c_K^{\mathrm{band}}\) or the band ratios \emph{as evidence} for or
against those conjectures is forbidden. They may be computed as context. What
remains permitted, and useful, are measurements whose answer differs between
competing mechanisms \emph{in a regular flow} --- questions about a mechanism,
not proxies for blow-up.

\begin{obsv}[What Layer~4 is still good for]
\label{obs:layer4-value}
Wall~6 indicts a family of measurements, not the layer. The diagnostic suite
itself is \OPEN{} and valid for other purposes, and several of its results are
of the mechanism kind Wall~6 permits: the \(\varepsilon\)-independence of
\(\delta\), which is a fact about the equations; the sampling-bias discovery;
the resolution correction to \(c_K\); the early-transient explanation of the
band spikes; the demonstration that free viscous decay does not suppress
gradient concentration. Each of those told the programme something true. What
none of them could do was decide a conditional statement about a regime no
regular solver can enter.
\end{obsv}
